% Owned proposed fragment B.129.
\subsection*{An invariant class with no fixed lift in the actual closed kernel}

\statement{B.129}{Proved}{A residue-three free-fiber quotient gives an actual equivariant free-two retract and proves nonlinear nonlifting in both categories}
\label{stmt:B129}
Use \(F,L,N,W,J,T\) and \(\delta\) from B.123 and B.127.
Write
\[
 \overline L=\ker(\widehat F\to\widehat T),\qquad
 \overline N=\operatorname{cl}_{\widehat F}(N),\qquad
 \overline J=\overline L/\overline N
             =\ker(\widehat W\to\widehat T).
\]
There are actual equivariant split epimorphisms
\begin{equation}\label{eq:B129-retracts}
 q:J\longrightarrow F(A,B),\qquad
 \widehat q:\overline J\longrightarrow\widehat{F(A,B)},
\end{equation}
where the target involution exchanges its two positive free
generators. Their marked sections send
\begin{equation}\label{eq:B129-section}
 A\longmapsto A_0=xax^{-1},\qquad
 B\longmapsto B_0=\bar a=cac .
\end{equation}
In particular the actual subgroup generated by \(A_0,B_0\)
has its full induced profinite topology in \(W\) and in the
ordinary one-neck group.

The classes \([A_0B_0]\) in \(J_{\rm ab}\) and in the actual
profinite abelianization
\[
 \overline J^{\,\rm ab}
   =\overline J/\overline{[\overline J,\overline J]}
\]
are invariant under \(\delta\), but have no representatives in
\(J^\delta\) and \(\overline J^\delta\), respectively. The
completed assertion uses the actual closed kernel, without
identifying it with \(\widehat J\) or identifying its abelianization
with a completion of an infinite discrete module. It constructs
no fixed element outside the source image.

\begin{proof}
\emph{B.129.1. The uniform color and level invariant.}
Let \(\chi:T\to\mathbb Z\) be given by
\(\chi(f)=\chi(g)=0\), \(\chi(x)=1\). The complete signed source
maps of B.123 and their retained heads in B.127 satisfy
\begin{equation}\label{eq:B129-level-invariant}
 \chi\phi_\sigma=\chi,\qquad
 \chi(h_{\sigma,i})=0
 \quad(\sigma=m,m^{-1},v,v^{-1};\ i\in\mathcal U).
\end{equation}
The first equations follow on the three retained generators;
the second hold for every entry of the exact head table.
The nonlinear formula in B.127 therefore shows that every
signed source word preserves the positive peripheral color
and its level \(\chi(t)\). Conjugation by a retained word
adds the same integer to both levels of a node, and conjugation
by \(L\) changes only its nonlinear connectors.

The complete forty-node input has only four pairs whose two
colors are \(a,\bar a\). Reordering the two endpoints so that
the \(a\)-color comes first, their exact levels are
\[
\begin{array}{c|c|c|c}
 \text{regional node}&\text{ordinary sheet}&
    (\chi(t_a),\chi(t_{\bar a}))&
       \chi(t_a)-\chi(t_{\bar a})\\ \hline
 -,\ \text{diagonal }a&0&(0,0)&0\\
 -,\ \text{diagonal }a&1&(0,0)&0\\
 +,\ \text{diagonal }a&0&(0,1)&-1\\
 +,\ \text{diagonal }a&1&(-1,0)&-1
\end{array}
\]
Every other node has a factor of another color. Thus the
relevant difference is always \(0\) or \(-1\), for the entire
normal closure including all signed source and retained-deck
translates. The connectors in B.127 remain arbitrary actual
words; they have not been removed.

\emph{B.129.2. A finite-rank continuous construction.}
Let \(F_{24}\) be the free group on
\[
 y_{i,r},\qquad i\in\mathcal U,\quad r\in\mathbb Z/3,
\]
and let \(s\), of order three, act by
\(sy_{i,r}s^{-1}=y_{i,r+1}\). There is a homomorphism
\begin{equation}\label{eq:B129-Phi}
 \Phi:F\longrightarrow F_{24}\rtimes C_3(s),\qquad
 U_i\mapsto y_{i,0},\quad x\mapsto s,\quad f,g\mapsto1.
\end{equation}
Its cyclic coordinate is \(\chi\) modulo three, so \(\Phi(L)\)
lies in the normal free fiber \(F_{24}\). If
\(\Phi(w)=ns^r\), then
\[
 \Phi(wU_i^\epsilon w^{-1})
     =n y_{i,r}^{\epsilon}n^{-1},\qquad
 r=\chi(\pi(w))\pmod3.
\]
This retains the entire nonlinear conjugator \(n\).

The universal property gives a continuous map
\[
 \widehat\Phi:\widehat F\longrightarrow
                \widehat F_{24}\rtimes C_3.
\]
Its restriction to \(\overline L\) lands in the closed free
fiber \(\widehat F_{24}\). No continuity of a projection onto
one individual coordinate of the infinite Schreier basis is
being assumed.

\emph{B.129.3. Exact involution and projection.}
Define
\begin{equation}\label{eq:B129-Delta}
 \Delta(s)=s^{-1},\qquad
 \Delta(y_{i,r})=y_{\bar i,\,1-r}.
\end{equation}
It squares to the identity and respects the shift relation:
both images of \(sy_{i,r}s^{-1}=y_{i,r+1}\) are
\(y_{\bar i,-r}\). On all eleven original generators
\(\Delta\Phi=\Phi\delta\); for example the image of
\(\delta(U_i)=xU_{\bar i}x^{-1}\) is \(y_{\bar i,1}\).
The same identities hold continuously.

Let \(\rho:F_{24}\to F(A,B)\) keep precisely
\[
 y_{a,1}\mapsto A,\qquad y_{\bar a,0}\mapsto B
\]
and kill the other twenty-two generators. These two retained
letters form a whole \(\Delta\)-orbit, so \(\rho\) and
\(\widehat\rho\) intertwine \(\Delta\) with the positive basis
swap. This projection is used on the free fiber only. It is
not asserted to extend across the \(C_3\)-action.

\emph{B.129.4. The entire actual relation subgroup is killed.}
Consider any defining node with an arbitrary signed source
translate and arbitrary ambient conjugator, using the exact
normal-closure description of B.127. It is a commutator of
two conjugated family cores. If one color differs from both
\(a,\bar a\), that core is killed by \(\rho\). If both colors
are relevant, retaining both cores would require the
\(a\)-level to be \(1\) and the \(\bar a\)-level to be \(0\)
modulo three. Their difference would be \(1\), whereas
B.129.1 gives only \(0\) or \(2\) modulo three.

Thus at least one core of every actual translated node is
killed. Its conjugate is killed regardless of its nonlinear
connector, so the complete commutator vanishes. Consequently
\(\rho\Phi|_L\) kills \(N\), and the closed kernel of
\(\widehat\rho\widehat\Phi|_{\overline L}\) kills
\(\overline N\). This gives both maps in
\eqref{eq:B129-retracts} with their actual quotient topology.
The proof is uniform over every source and deck word; no
finite orbit cutoff is used, and the source operators need
not descend to the auxiliary semidirect target.

\emph{B.129.5. Literal sections and full induced topology.}
The two words in \eqref{eq:B129-section} map exactly to \(A,B\),
and satisfy \(\delta(A_0)=B_0\), \(\delta(B_0)=A_0\) already
as free-eleven words. They therefore define equivariant
sections. In the profinite case the universal property of
\(\widehat F_2\) gives a continuous section into the actual
\(\overline J\), whose composite with \(\widehat q\) is the
identity on both free generators.

There is also a direct open-subgroup interpretation. Put
\[
 F^{(3)}=\ker(F\xrightarrow{\chi\pi}\mathbb Z/3).
\]
The map \(\rho\Phi\) is defined on all of \(F^{(3)}\), kills
\(N\), and descends to \(W^{(3)}=F^{(3)}/N\), an actual
index-three subgroup of \(W\). It retracts \(W^{(3)}\) onto
the same actual free-two subgroup. Its full completion is
an open subgroup of \(\widehat W\); restricting its completed
retraction to \(\overline J\) recovers \(\widehat q\).
Together with B.123's full topology on \(W\), this proves the
asserted induced topology in the whole ordinary group.

\emph{B.129.6. The target fixed subgroup and nonlifting.}
For the involution \(\epsilon\) swapping \(A,B\), elimination
of \(B=\epsilon A\epsilon\) gives the exact presentations
\[
 F_2\rtimes\langle\epsilon\rangle
       =C_2(\epsilon)*\mathbb Z(A),\qquad
 \widehat F_2\rtimes\langle\epsilon\rangle
       =C_2(\epsilon)\amalg\widehat{\mathbb Z}(A).
\]
The second equality follows by the universal profinite
presentations; its orientation kernel has the two indicated
free generators. The discrete free-product normal form
shows that the centralizer of \(\epsilon\) is
\(\langle\epsilon\rangle\). In the full free profinite product
the same factor-centralizer equality is Theorem B' of
\citet{A_HR1985}, applied to the nontrivial element of its
\(C_2\) factor. Its intersection with the orientation kernel
is trivial. Therefore
\[
 F_2^{\rm swap}=1,\qquad
 \widehat F_2^{\,\rm swap}=1,\qquad
 q(J^\delta)=1,\quad
 \widehat q(\overline J^\delta)=1.
\]

The class of \(A_0B_0\) is invariant in either actual
abelianization because \(\delta\) exchanges its two factors.
Under the abelianized retraction it maps to \((1,1)\) in
\(\mathbb Z^2\), respectively \(\widehat{\mathbb Z}^2\).
A fixed representative would map to the identity in the
free-two target and hence to zero in its abelianization,
a contradiction. This proves both nonlifting assertions.

Since the actual source image is fixed, its image under the
equivariant free-two quotient is trivial. The argument does
not construct a fixed element outside that source, and is
compatible with either possible outcome of the open equalities
in B.123.
\end{proof}

\begingroup\raggedright\footnotesize
Complete proof and exact residue-three checks:
\path{work_CE/compact_fixed_class_nonlift26.txt},
\path{work_CE/code/check_compact_nonlift_residue3_26.py},
\path{work_CE/code/check_compact_nonlift_residue3_26.out},
\path{work_CE/next_symmetric_nodes/compact_nonlift_residue3_26.json}.
The finite certificate retains all nonlinear connectors and
checks all forty nodes in each of three common residues.
The infinite signed-orbit implication is the proof in
B.129.1 and B.129.4, not an inference from those finite checks.
\par\endgroup
